\documentclass{article}

\usepackage{amsmath,amssymb}
\usepackage{xurl}
\usepackage[hidelinks]{hyperref}
\setlength{\emergencystretch}{2em}

% Shared commands for notes in docs/paper-gaps/.

\DeclareUnicodeCharacter{2080}{\ifmmode{}_{0}\else\textsubscript{0}\fi}
\DeclareUnicodeCharacter{2081}{\ifmmode{}_{1}\else\textsubscript{1}\fi}
\DeclareUnicodeCharacter{2082}{\ifmmode{}_{2}\else\textsubscript{2}\fi}
\DeclareUnicodeCharacter{2099}{\ifmmode{}_{n}\else\textsubscript{n}\fi}

\newcommand{\Lean}{\textsc{Lean}}
\ExplSyntaxOn
\NewDocumentCommand{\leanid}{m}{
  \ifmmode
    \textsf{\detokenize{#1}}
  \else
    \group_begin:
    \urlstyle{sf}
    \tl_set:Nn \l_tmpa_tl {#1}
    \tl_replace_all:Nnn \l_tmpa_tl {\_} {_}
    \exp_args:NV \path \l_tmpa_tl
    \group_end:
  \fi
}
\ExplSyntaxOff
\providecommand{\tr}{\operatorname{tr}}
\newcommand{\tnleanIssueBase}{https://github.com/LionSR/TNLean/issues/}
\newcommand{\tnleanPrBase}{https://github.com/LionSR/TNLean/pull/}
\newcommand{\tnleanDocBase}{https://sirui-lu.com/TNLean/docs/}
\newcommand{\ghissue}[1]{\href{\tnleanIssueBase#1}{GitHub issue \##1}}
\newcommand{\ghpr}[1]{\href{\tnleanPrBase#1}{PR \##1}}
\newcommand{\doclink}[1]{\href{\tnleanDocBase#1.html}{\path{#1}}}

% Dirac notation and the identity operator, as in blueprint/src/macros/common.tex.
% \providecommand keeps note-local definitions authoritative.
\usepackage{dsfont}
\providecommand{\ket}[1]{|#1\rangle}
\providecommand{\bra}[1]{\langle #1|}
\providecommand{\braket}[2]{\langle #1 | #2 \rangle}
\providecommand{\ketbra}[2]{|#1\rangle\!\langle #2|}
\providecommand{\Id}{\mathds{1}}
\providecommand{\one}{\mathds{1}}
\providecommand{\C}{\mathbb{C}}
\providecommand{\MN}[1]{M_{#1}(\C)}
\providecommand{\MD}{M_D(\C)}

% Verdict marker: \gapnote{<kind>}{<status>}. Kind is the policy
% classification (clarification, local-correction, scope-restriction,
% unfaithful, false-source, open-gap); status is open, wip (elimination
% underway), resolved, or historical. Machine-read by the site index;
% prints a verdict line.
\newcommand{\gapnote}[2]{\par\noindent\textbf{Verdict:} #1 (#2).\par}


% Fallbacks keep this note compilable when a different command.tex is found first.
\providecommand{\leanid}[1]{\textsf{\detokenize{#1}}}
\providecommand{\gapnote}[2]{\par\noindent\textbf{Verdict:} #1 (#2).\par}

\title{The Bond Dimension of the Frustration-Free Matrix Product State Approximation of Arad, Kitaev, Landau and Vazirani}
\author{}
\date{2026-11-04}

\begin{document}
\maketitle
\gapnote{local-correction}{resolved}

\section{The source statement}

Arad, Kitaev, Landau and Vazirani~\cite{Arad2013AreaLaw},
arXiv:1301.1162v1, \path{1dfrustrated34.tex}, Section~5, lines~851--870,
consider a frustration-free nearest-neighbour chain of $n$ sites of
dimension $d$, with terms $0\le H_i\le1$, a unique ground state $\Gamma$ and
gap $\varepsilon$. They state that $\Gamma$ is within $1/\mathrm{poly}(n)$ of a
matrix product state of bond dimension
\begin{align}
    \exp\bigl(O(\varepsilon^{-1/3}\log^{2/3}n)\bigr).
    \label{eq:aklv13_ffmps_printed}
\end{align}
The argument (lines~862--870) starts from a state of constant overlap with
$\Gamma$ and entanglement rank $\exp(O(\log^3d/\varepsilon))$, and applies an
approximate ground-state projector with shrinking factor
$\Delta=1/\mathrm{poly}(n)$. It takes $\ell=s^2$ and states that Lemma~4.2
of Ref.~\cite{Arad2013AreaLaw} gives $\Delta=e^{-c\varepsilon^{-1/2}s^{3/2}}$ and rank
$D=\exp(s\log d)$; then $s=O(\varepsilon^{-1/3}\log^{2/3}n)$.

\section{The exponent of the gap at line 867}

The exponent of $\varepsilon$ in $\Delta$ is $+1/2$; the sign at line~867 is
a misprint. The Chebyshev polynomial of degree $\ell$ of the clipped
Hamiltonian, whose norm is at most $s+2t$ with $t=O(1/\varepsilon)$, gives
$\Delta=\exp(-\ell\sqrt{\varepsilon}/\sqrt{s+2t})$ (line~831). For $\ell=s^2$
this is $\exp(-c\,\varepsilon^{1/2}s^{3/2})$ with $c=\sqrt{s}/\sqrt{s+2t}$, as
in the proof of Theorem~4.3 of Ref.~\cite{Arad2013AreaLaw} (line~835). The choice of $s$ at
line~868 makes $\varepsilon^{1/2}s^{3/2}=\log n$, which is the balance
needed for $\Delta=1/\mathrm{poly}(n)$.

The constant $c$ is bounded below only when $s$ is at least of order
$1/\varepsilon$. With $s$ of order $\varepsilon^{-1/3}\log^{2/3}n$ this needs
$\log n$ of order at least $1/\varepsilon$. The starting rank
$\exp(O(\log^3d/\varepsilon))$ is also omitted from
(\ref{eq:aklv13_ffmps_printed}). Both are covered by an additive term in
$\log B$ that is independent of $n$.

\section{The rank of the second projector}

Lemma~4.2 of Ref.~\cite{Arad2013AreaLaw} (line~590) bounds the entanglement rank of a polynomial
of degree $\ell$ by $(d\ell)^{O(\max(\ell/s,\sqrt{\ell}))}$. At $\ell=s^2$ this
is $\exp(O(s(\log d+\log s)))$, and the proof of Theorem~4.3 keeps
the same factor $\log(s^2+s)$ (lines~840--841). The count is sharp for the
method: the binomial coefficient $\binom{s^2+s}{s}$ in the proof is at least
$s^s$. Line~868 writes $D=\exp(s\log d)$ and drops $\log s$. With
$s=\varepsilon^{-1/3}\log^{2/3}n$ the factor $\log s$ is of order
$\log\log n+\log(1/\varepsilon)$. The source's lemmas therefore give
\begin{align}
    \log B=O\bigl(\varepsilon^{-1/3}\log^{2/3}n\,
        (\log d+\log\log n+\log(1/\varepsilon))\bigr),
    \label{eq:aklv13_ffmps_source}
\end{align}
not (\ref{eq:aklv13_ffmps_printed}). This is the factor hidden by the tilde
in the general bound $e^{\tilde O(\log^{3/4}n/\varepsilon^{1/4})}$ of the same
paper (lines~176 and 1503--1517).

\section{What is formalized}

For a nearest-neighbour chain of $L$ sites of dimension $q$ with terms
$0\le h_j\le1$, a unit vector $\Omega$ with $H\Omega=0$ for $H=\sum_jh_j$,
$H\ge g(I-|\Omega\rangle\langle\Omega|)$ with $g>0$ and $g$ in the spectrum of
$H$, and for $0<\epsilon\le1$, put
\begin{align}
    u=1+\log q+\log\frac1g,\qquad
    z=\log\frac{2048(L+1)}{\epsilon^2}.
    \label{eq:aklv13_ffmps_parameters}
\end{align}
There are a universal constant $C>0$ and a bond dimension $D$ with
\begin{align}
    \log D\le C\Bigl[\frac{u^3}{g}
        +\Bigl(\frac{z^2}{g}\Bigr)^{1/3}\bigl(u+\log(z+1)\bigr)\Bigr]
    \label{eq:aklv13_ffmps_formal}
\end{align}
and a normalized open-boundary matrix product state $\Phi$ of bond dimension
at most $D$ with $\lVert\Phi-\Omega\rVert\le\epsilon$.\footnote{%
    \leanid{TNLean.MPS.AreaLaw.exists_aklv_frustrationFree_sequential_mps_approximation}
    in \doclink{TNLean/MPS/AreaLaw/AKLV/FrustrationFreeMPSApproximation}.}
The hypotheses are those of the source: $g$ is the gap itself, and $q\ge2$
and $g\le1$ are derived.

The proof follows lines~862--868. At each cut a product vector with squared
overlap at least $1/(2D_0)$ with $\Omega$ is filtered by the projector of the
entropy theorem, which gives a vector within squared distance $1/16$ of
$\Omega$ and of rank at most $D_0$, with $\log D_0\le Ru^3/g$. A second
projector of degree $\ell=s^2$ on a window of width
$s=s_0+2\lceil6(z^2/g)^{1/3}\rceil$, with $s_0\ge1/g$ the window of the entropy
theorem, has shrinking factor at most $\epsilon^2/(1024(L+1))$ at the clipping
height $1024/g$. Here the exponent $+1/2$ of the gap is used. If the window
does not fit, $\Omega$ itself has rank at most $q^{s-\lfloor s/2\rfloor}$.
Sequential truncation of the leading Schmidt vectors then gives $\Phi$.

\section{Comparison}

Take $\epsilon=n^{-c}$ with $n=L$. Then $z=O(\log n)$, and for fixed $q$ and
$g$ the bound (\ref{eq:aklv13_ffmps_formal}) reads
\begin{align}
    \log D=O\bigl(g^{-1/3}(\log n)^{2/3}\log\log n\bigr).
    \label{eq:aklv13_ffmps_asymptotic}
\end{align}
This is (\ref{eq:aklv13_ffmps_source}) with the factor $\log s$ written out.
It is smaller than the general bound $(\log n)^{3/4}$ of the same paper. The
printed bound (\ref{eq:aklv13_ffmps_printed}), without the factor
$\log\log n$, is not formalized and is not implied by Lemma~4.2 of Ref.~\cite{Arad2013AreaLaw}.
The gap enters (\ref{eq:aklv13_ffmps_formal}) also through $\log(1/g)$ in
$u$ and through the term $u^3/g$, which are independent of $n$; the inverse-gap
logarithms are retained as in the entropy theorem,
\path{docs/paper-gaps/aklv13_two_block_robustness.tex}.

\section{Verdict}

A local correction, resolved: the exponent at line~867 is read as $+1/2$, and
the bond dimension of Section~5 is formalized with the factor $\log\ell$ of
Lemma~4.2 of Ref.~\cite{Arad2013AreaLaw}, under the source's hypotheses.

\bibliographystyle{plain}
\bibliography{references}

\end{document}
